\documentclass[11pt]{article}
\usepackage[a4paper,margin=1in]{geometry}
\usepackage[T1]{fontenc}
\usepackage[utf8]{inputenc}
\usepackage{setspace}
\usepackage{hyperref}

\setstretch{1.1}
\setlength{\parskip}{0.6em}
\setlength{\parindent}{0pt}

\begin{document}

\begin{flushright}
April 29, 2026
\end{flushright}

Editors \\
\textit{PLOS ONE}

Dear Editors,

Please consider our manuscript entitled

\textbf{``Visual anchoring predicts multimodal LLM annotation reliability: A diagnostic study for classroom behavior recognition''}

for publication in \textit{PLOS ONE}.

This manuscript presents a systematic diagnostic study of multimodal large
language model (MLLM) pseudo-labeling for classroom behavior recognition. Using
three publicly available SCB sub-datasets, we evaluate two representative
7B-scale MLLMs, multiple pseudo-label filtering strategies, protocol-matched
zero-shot CLIP baselines, and 18 CLIP ViT-L/14 fine-tuning conditions covering
linear probing and LoRA. The study is designed to answer a practical and timely
question: when can MLLMs be used as low-cost annotators for classroom behavior
recognition, and when do they fail despite apparent visual-language capability?

The main finding is that pseudo-label reliability is strongly associated with
the visual anchoring of each behavior category. Visually explicit classes such
as blackboard-writing and teacher can be labeled reliably, whereas
intent-dependent classroom interactions such as answering, guiding, and
on-stage interaction remain difficult for both tested MLLMs. We further show
that agreement-based filtering is not a universally reliable denoising strategy,
because it can sharply reduce retained data volume and still produce biased
retained subsets. Downstream experiments demonstrate dataset-dependent gains
from pseudo-label fine-tuning relative to zero-shot CLIP, while also showing a
large remaining gap between pseudo-label and ground-truth training for
intent-dependent teacher behaviors.

We believe the manuscript is well suited to \textit{PLOS ONE} because it reports
a technically sound, reproducible empirical evaluation with clear implications
for educational AI, computer vision annotation, and responsible use of
general-purpose MLLMs. The contribution is not a claim of a new model, but a
careful characterization of where an emerging annotation paradigm works, where
it fails, and how these limits can be diagnosed before large-scale labeling is
undertaken.

This submission is original, has not been published previously, and is not under
consideration elsewhere. All authors have approved the submitted version and
agree to its submission to \textit{PLOS ONE}. The authors declare no competing
interests.

Funding statement: This research was funded by the Education Department of
Hunan Province, grant number HNJG-2022-0608.

Data and code availability: The SCB dataset is publicly available at
\url{https://huggingface.co/datasets/wintonYF/SCB-Dataset}. The experiment code
is available at \url{https://github.com/zhanglizhuo/llm_annotation}.

Thank you for your time and consideration.

Sincerely,

Yan Ma and Lizhuo Zhang

Corresponding author:
Lizhuo Zhang \\
School of Information and Intelligent Science and Technology, \\
Hunan Agricultural University, Changsha 410128, China \\
Email: zhanglizhuo@hunau.edu.cn

\end{document}